\section{2017 Analysis \& Differential Equations}

\subsection{Individual}

\begin{exercise}
\label{analysis:2017:1}
Suppose $f$ is integrable on $[-\pi,\pi]$. Prove that 
\[
\sum_{n=-\infty}^{\infty} a_n r^{|n|}e^{inx}
\]
tends to $f(x)$ for a.e. $x$, as $r \to 1$, $r < 1$. Here $a_n = \frac{1}{2\pi}\int_{-\pi}^{\pi}f(x)e^{-inx}\,dx$.
\end{exercise}

\begin{solution}
\textbf{Prerequisites:}
To understand this problem, you should be familiar with:
\begin{itemize}
    \item \textbf{Fourier Series:} Representing an integrable function as a sum of complex exponentials.
    \item \textbf{Poisson Kernel:} A specific kernel $P_r(\theta)$ used to define the Abel sum of a series.
    \item \textbf{Approximate Identities:} Sequences of functions that "concentrate" their mass at zero, acting like a Dirac delta function in the limit.
    \item \textbf{Lebesgue Points:} Points where a function is "well-behaved" in an average sense; almost every point for an integrable function is a Lebesgue point.
\end{itemize}

\textbf{Steps:}
\begin{enumerate}
    \item \textbf{Identify the Abel Sum:} The given sum $S_r(x) = \sum_{n=-\infty}^{\infty} a_n r^{|n|}e^{inx}$ is precisely the Abel sum of the Fourier series of $f$. We replace the coefficients $a_n$ with their integral definition: $a_n = \frac{1}{2\pi}\int_{-\pi}^{\pi}f(t)e^{-int}\,dt$.
    \item \textbf{Express as a Convolution:} By substituting the definition of $a_n$ into the sum and swapping the order of summation and integration, we get:
    \[ S_r(x) = \frac{1}{2\pi} \int_{-\pi}^{\pi} f(t) \left( \sum_{n=-\infty}^{\infty} r^{|n|} e^{in(x-t)} \right) dt \]
    The term in the parentheses is the \textbf{Poisson Kernel}, $P_r(x-t)$.
    \item \textbf{Analyze the Poisson Kernel:} The Poisson kernel integrates to 1 and satisfies the properties of an approximate identity as $r \to 1^-$. 
    \item \textbf{Apply Convergence Results:} Because $P_r$ is an approximate identity, $(f * P_r)(x)$ converges to $f(x)$ at every \textbf{Lebesgue point} of $f$. 
    \item \textbf{Conclusion:} By the Lebesgue Differentiation Theorem, almost every point $x$ is a Lebesgue point for $f \in L^1$. Thus, $S_r(x) \to f(x)$ for a.e. $x$.
\end{enumerate}

\textbf{Intuition:}
Multiplying Fourier coefficients by $r^{|n|}$ "dampens" high-frequency oscillations. As $r \to 1$, we slowly re-introduce these frequencies. Since the Poisson kernel is a positive averaging operator that focuses on the value at $x$, it successfully reconstructs the function almost everywhere.
\end{solution}

\begin{exercise}
\label{analysis:2017:2}
Let $H$ be a Hilbert space equipped with an inner product $(\cdot,\cdot)$ and a norm $\|\cdot\| = (\cdot,\cdot)^{1/2}$. A sequence $\{f_k\}$ converges to $f \in H$ if $\|f_k-f\| \to 0$. A sequence $\{f_k\} \subset H$ is said to converge weakly to $f \in H$ if $(f_k,g) \to (f,g)$ for any $g \in H$. Prove the following statements:
\begin{enumerate}
    \item $\{f_k\}$ converges to $f$ if and only if $\|f_k\| \to \|f\|$ and $\{f_k\}$ converges weakly to $f$.
    \item If $H$ is a finite dimensional Hilbert space, then the weak convergence implies convergence. Give a counterexample to show that weak convergence does not necessarily imply convergence in an infinite dimensional Hilbert space.
\end{enumerate}
\end{exercise}

\begin{solution}
\textbf{Prerequisites:}
\begin{itemize}
    \item \textbf{Strong vs Weak Convergence:} Strong convergence is distance-based; weak convergence is projection-based.
    \item \textbf{Inner Product Properties:} $\|u-v\|^2 = \|u\|^2 - 2\text{Re}(u, v) + \|v\|^2$.
    \item \textbf{Orthonormal Basis:} In infinite dimensions, basis vectors stay far apart but their projections onto any fixed vector go to zero.
\end{itemize}

\textbf{Steps:}
\begin{enumerate}
    \item \textbf{Part 1 ($\implies$):} If $f_k \to f$ strongly, then $(f_k, g) \to (f, g)$ by Cauchy-Schwarz, and $\|f_k\| \to \|f\|$ by the triangle inequality.
    \item \textbf{Part 1 ($\impliedby$):} Expand $\|f_k - f\|^2 = \|f_k\|^2 - (f_k, f) - (f, f_k) + \|f\|^2$. As $k \to \infty$, this becomes $\|f\|^2 - \|f\|^2 - \|f\|^2 + \|f\|^2 = 0$.
    \item \textbf{Part 2 (Finite Dim):} Weak convergence implies coordinate-wise convergence. In finite dimensions, this is equivalent to norm convergence.
    \item \textbf{Part 2 (Infinite Dim Counterexample):} In $\ell^2$, the unit basis vectors $e_k$ converge weakly to 0 (since every component of a square-summable vector eventually vanishes), but $\|e_k\| = 1 \not\to 0$.
\end{enumerate}

\textbf{Intuition:}
Strong convergence means the vectors are close. Weak convergence means they "look like" the limit from any fixed direction. In infinite dimensions, a sequence can "run away" into new dimensions (like $e_k$) so that it looks like 0 from every old direction, while its length stays 1.
\end{solution}

\begin{exercise}
\label{analysis:2017:3}
Let $f:\mathbb{C}\setminus\{0\}\to\mathbb{C}$ be a holomorphic function and
\[
|f(z)| \leq |z|^2 + \frac{1}{|z|^{1/2}},
\]
for $z$ near 0. Determine all such functions.
\end{exercise}

\begin{solution}
\textbf{Prerequisites:}
\begin{itemize}
    \item \textbf{Riemann's Removable Singularity Theorem:} A holomorphic function that is bounded near a point has a removable singularity there.
    \item \textbf{Laurent Coefficients:} The negative Laurent coefficients can be estimated by Cauchy's integral formula on small circles.
    \item \textbf{Entire Function:} A function holomorphic on the whole complex plane.
\end{itemize}

\textbf{Steps:}
\begin{enumerate}
    \item \textbf{Estimate the principal part:} Write the Laurent expansion near $0$ as $f(z)=\sum_{m=-\infty}^{\infty}a_mz^m$. For $k\geq 1$, Cauchy's formula on $|z|=r$ gives
    \[
    |a_{-k}|
    \leq r^k\max_{|z|=r}|f(z)|
    \leq r^k\left(r^2+r^{-1/2}\right)
    =r^{k+2}+r^{k-1/2}.
    \]
    Letting $r\to 0$ gives $a_{-k}=0$ for every $k\geq 1$.
    \item \textbf{Remove the singularity:} Since all negative Laurent coefficients vanish, the singularity at $0$ is removable. Thus $f$ extends holomorphically to all of $\mathbb{C}$.
    \item \textbf{Identify all possibilities:} Conversely, any entire function is bounded in a sufficiently small disk around $0$, while $|z|^2+|z|^{-1/2}\to\infty$ as $z\to0$. Hence every entire function, restricted to $\mathbb{C}\setminus\{0\}$, satisfies the stated local bound near $0$.
\end{enumerate}

\textbf{Intuition:}
Complex functions are very rigid. If a function doesn't blow up fast enough to be a pole (like $1/z$), it must not be blowing up at all. A $1/\sqrt{z}$ growth is too "tame" for a true singularity.
\end{solution}

\begin{exercise}
\label{analysis:2017:4}
Find a conformal mapping which maps the region $\{z\mid |z-i|<\sqrt{2}, |z+i|<\sqrt{2}\}$ onto the unit disk.
\end{exercise}

\begin{solution}
\textbf{Prerequisites:}
\begin{itemize}
    \item \textbf{Möbius Transformations:} Maps of the form $(az+b)/(cz+d)$.
    \item \textbf{Conformal Geography:} Mapping lenses to sectors, sectors to half-planes, and half-planes to disks.
\end{itemize}

\textbf{Steps:}
\begin{enumerate}
    \item \textbf{Find intersection points:} The circles $|z-i|=\sqrt{2}$ and $|z+i|=\sqrt{2}$ intersect at $z=1$ and $z=-1$.
    \item \textbf{Map to a sector:} The map $w_1 = (z-1)/(z+1)$ sends $1 \to 0$ and $-1 \to \infty$. The "lens" becomes a wedge of angle $\pi/2$ (since the circles are orthogonal).
    \item \textbf{Normalize and Square:} Rotate the wedge to the positive real axis $w_2 = -w_1$ and square it $w_3 = w_2^2$ to get a half-plane.
    \item \textbf{Final map:} Map the right half-plane to the disk. The composite map simplifies to $f(z) = -2z/(1+z^2)$.
\end{enumerate}

\textbf{Intuition:}
We first "unbend" the circular arcs into straight lines by moving the corners to 0 and $\infty$. This turns the lens into a pizza slice. We then double the angle of the slice until it fills a half-plane, and finally "fold" that infinite plane into the unit circle.
\end{solution}

\begin{exercise}
\label{analysis:2017:5}
If $E$ is a compact set in a region $\Omega$, prove that there exists a constant $M>0$, depending only on $E$ and $\Omega$, such that every positive harmonic function $u(z)$ in $\Omega$ satisfies $u(z_2)\leq Mu(z_1)$ for any two points $z_1,z_2\in E$.
\end{exercise}

\begin{solution}
\textbf{Prerequisites:}
\begin{itemize}
    \item \textbf{Harnack's Inequality:} For a disk $B(a, R)$, $u(z) \leq \frac{R+r}{R-r} u(a)$.
    \item \textbf{Compactness:} Finite subcovers exist for any open cover.
\end{itemize}

\textbf{Steps:}
\begin{enumerate}
    \item \textbf{Connectedness convention:} We use the standard complex-analysis convention that a region is connected. This is needed. If $\Omega$ were allowed to be disconnected and $E$ met two different components, a positive harmonic function could be multiplied by an arbitrary positive constant on one component, so no uniform comparison between the two components would be possible.
    \item \textbf{Local bound:} Around any point in $E$, pick a small disk where $u(z) \leq C u(a)$.
    \item \textbf{Covering:} Cover the compact set $E$ with a finite number of such overlapping disks.
    \item \textbf{Chaining:} Since $\Omega$ is connected and $E$ is compactly contained in $\Omega$, finitely many Harnack balls can be chosen so that any two points of $E$ are connected by a chain of at most $N$ balls staying a positive distance from $\partial\Omega$. The ratio $u(z_2)/u(z_1)$ is bounded by the product of the Harnack constants from each step in the chain.
\end{enumerate}

\textbf{Intuition:}
Harmonic functions are very smooth. If the function is positive, it can't drop to 0 or spike to infinity too quickly. Over a compact set, you stay "safely away" from the boundary where things could go wrong, so the values at any two points must be comparable.
\end{solution}

\begin{exercise}
\label{analysis:2017:6}
\begin{enumerate}
\item For any bounded domain $\Omega\subset\mathbb{R}^n$, prove that there exists a smallest constant $C(\Omega)$ such that
\[
\int_{\Omega}|u|^2\,dx
\leq C(\Omega)\int_{\Omega}\sum_{i=1}^n
\left|\frac{\partial u}{\partial x_i}\right|^2\,dx
\]
for every $u\in H_0^1(\Omega)$.
\item Let
\[
\Pi=\{(x,y):0<x<a,\ 0<y<b\}.
\]
Show that
\[
C(\Pi)\geq \frac{a^2b^2}{\pi^2(a^2+b^2)}.
\]
\end{enumerate}
\end{exercise}

\begin{solution}
\textbf{Prerequisites:}
\begin{itemize}
    \item \textbf{Rayleigh Quotient:} $\lambda_1 = \inf \int |\nabla u|^2 / \int u^2$.
    \item \textbf{Dirichlet Eigenvalues:} The smallest $\lambda$ for $-\Delta u = \lambda u$ with $u=0$ on $\partial \Omega$.
\end{itemize}

\textbf{Steps:}
\begin{enumerate}
    \item \textbf{Define the sharp constant:} The Poincar\'e inequality on a bounded domain says that some finite constant works. The sharp one is
    \[
    C(\Omega)
    =
    \sup_{u\in H_0^1(\Omega),\,u\neq0}
    \frac{\int_\Omega |u|^2\,dx}{\int_\Omega |\nabla u|^2\,dx}.
    \]
    This is finite by the usual Poincar\'e inequality and is automatically the smallest admissible constant: any valid constant must dominate every quotient in the supremum, while the supremum itself gives the inequality.
    \item \textbf{Spectral interpretation:} Equivalently,
    \[
    C(\Omega)=\frac{1}{\lambda_1(\Omega)},\qquad
    \lambda_1(\Omega)=
    \inf_{u\in H_0^1(\Omega),\,u\neq0}
    \frac{\int_\Omega |\nabla u|^2\,dx}{\int_\Omega |u|^2\,dx}.
    \]
    Here $\lambda_1(\Omega)$ is the first Dirichlet eigenvalue of $-\Delta$.
    \item \textbf{Rectangle case:} For $\Pi=(0,a)\times(0,b)$, take the test function
    \[
    u(x,y)=\sin\frac{\pi x}{a}\sin\frac{\pi y}{b}.
    \]
    Then $u\in H_0^1(\Pi)$ and
    \[
    -\Delta u=\pi^2\left(\frac1{a^2}+\frac1{b^2}\right)u.
    \]
    Hence
    \[
    \frac{\int_\Pi |\nabla u|^2}{\int_\Pi |u|^2}
    =\pi^2\left(\frac1{a^2}+\frac1{b^2}\right).
    \]
    Therefore
    \[
    C(\Pi)
    \geq
    \frac{\int_\Pi |u|^2}{\int_\Pi |\nabla u|^2}
    =
    \frac{a^2b^2}{\pi^2(a^2+b^2)}.
    \]
    In fact this test function is the first eigenfunction, so equality holds.
\end{enumerate}

\textbf{Intuition:}
The constant $C(\Omega)$ measures how much a function can "bend" while staying zero at the boundary. A smaller domain forces more bending (higher gradient), making $C(\Omega)$ smaller.
\end{solution}

\subsection{Team}

\begin{exercise}
\label{analysis:2017:7}
Suppose that $f$ is an integrable function on the circle with 
\[
\hat{f}(n) = \frac{1}{2\pi} \int_0^{2\pi} f(x) e^{-inx} dx = 0,
\]
for all $n \in \mathbf{Z}$. Then $f(\theta) = 0$ whenever $f$ is continuous at the point $\theta$.
\end{exercise}

\begin{solution}
\textbf{Prerequisites:}
\begin{itemize}
    \item \textbf{Fejér's Theorem:} The Cesàro means of the Fourier series of a continuous function converge uniformly to the function.
    \item \textbf{Fourier Uniqueness:} If all coefficients are zero, the function is zero at its points of continuity.
\end{itemize}

\textbf{Steps:}
\begin{enumerate}
    \item \textbf{Relate to Fourier Series:} The condition $\hat{f}(n)=0$ means the entire Fourier series of $f$ is identically zero.
    \item \textbf{Use Summability:} The Fejér means $\sigma_N(f)(x)$, which are linear combinations of $\hat{f}(n)$, are all zero for any $x$.
    \item \textbf{Continuous Points:} At any point $\theta$ where $f$ is continuous, Fejér's theorem guarantees $\sigma_N(f)(\theta) \to f(\theta)$.
    \item \textbf{Conclusion:} Since the limit of a sequence of zeros is zero, $f(\theta) = 0$.
\end{enumerate}

\textbf{Intuition:}
If a function has no frequency components at all, it shouldn't exist. For integrable functions, this means it's zero almost everywhere, and specifically zero at every point where it doesn't have a jump or singularity.
\end{solution}

\begin{exercise}
\label{analysis:2017:8}
Let $F$ be continuous on $[a, b]$. Show that:
\begin{enumerate}
    \item Suppose $(D^+ F)(x) = \limsup_{h \to 0, h > 0} \frac{F(x+h) - F(x)}{h} > 0$ for every $x \in [a, b)$. Then $F$ is increasing on $[a, b]$.
    \item If $F'(x)$ exists for every $x \in (a, b)$ and $|F'(x)| \leq M$, then $|F(x) - F(y)| \leq M|x-y|$ and $F$ is absolutely continuous.
\end{enumerate}
\end{exercise}

\begin{solution}
\textbf{Prerequisites:}
\begin{itemize}
    \item \textbf{Mean Value Theorem:} If $f$ is differentiable, $(f(b)-f(a))/(b-a) = f'(c)$.
    \item \textbf{Dini Derivatives:} Generalizations of the derivative for non-differentiable functions.
    \item \textbf{Absolute Continuity:} A property stronger than continuity, ensuring the Fundamental Theorem of Calculus holds.
\end{itemize}

\textbf{Steps:}
\begin{enumerate}
    \item \textbf{Part 1 (Increasing):} Suppose $F(y) < F(x)$ for some $y > x$. Let $G(z) = F(z) + \epsilon z$ for small $\epsilon > 0$. The maximum of $G$ on $[x, y]$ cannot be at $y$, and the positive Dini derivative prevents it from being in the middle. This contradiction shows $F$ must be non-decreasing.
    \item \textbf{Part 2 (Lipschitz):} Apply the Mean Value Theorem to $F$ on any sub-interval $[x, y]$. Since $|F'(c)| \leq M$, we have $|F(y)-F(x)| \leq M|y-x|$.
    \item \textbf{Absolute Continuity:} Any Lipschitz continuous function is automatically absolutely continuous because the sum $\sum |F(x_i)-F(y_i)| \leq M \sum |x_i-y_i|$ can be made smaller than $\epsilon$ by choosing $\delta = \epsilon/M$.
\end{enumerate}

\textbf{Intuition:}
If the "slope" of a function is always pointing up (even in a weak sense like Dini derivatives), the function can never fall below its previous values. If the slope is bounded by $M$, the function can't grow or shrink faster than a line with slope $M$.
\end{solution}

\begin{exercise}
\label{analysis:2017:9}
Let $f$ be holomorphic on $|z|<1$ and assume $|f(z)|<1$ for $|z|<1$. If $\alpha,\beta\in\mathbb{C}$ satisfy $|\alpha|<1$, $|\beta|<1$, $\alpha\neq\beta$, and
\[
f(\alpha)=f(\beta)=0,
\]
prove that
\[
|f(z)|\leq
\left|\frac{z-\alpha}{1-\overline{\alpha}z}\right|
\left|\frac{z-\beta}{1-\overline{\beta}z}\right|,
\qquad |z|<1.
\]
\end{exercise}

\begin{solution}
\textbf{Prerequisites:}
\begin{itemize}
    \item \textbf{Schwarz Lemma:} If $f(0)=0$ and $|f| \leq 1$, then $|f(z)| \leq |z|$.
    \item \textbf{Blaschke Product:} A function that is 0 at specific points and has absolute value 1 on the boundary.
    \item \textbf{Maximum Modulus Principle:} A holomorphic function's maximum occurs on the boundary.
\end{itemize}

\textbf{Steps:}
\begin{enumerate}
    \item \textbf{First zero:} Let
    \[
    B_\alpha(z)=\frac{z-\alpha}{1-\overline{\alpha}z}.
    \]
    Since $f(\alpha)=0$, the quotient $h(z)=f(z)/B_\alpha(z)$ has a removable singularity at $\alpha$.
    \item \textbf{Bound the quotient:} The Schwarz--Pick lemma applied to $f$ and the point $\alpha$ gives
    \[
    |f(z)|\leq |B_\alpha(z)|.
    \]
    Hence $|h(z)|\leq1$ on the disk. Also $h(\beta)=0$, since $\beta\neq\alpha$ and $f(\beta)=0$.
    \item \textbf{Second zero:} Apply Schwarz--Pick again to the holomorphic self-map $h$ with $h(\beta)=0$. We get
    \[
    |h(z)|\leq
    \left|\frac{z-\beta}{1-\overline{\beta}z}\right|.
    \]
    Multiplying by $|B_\alpha(z)|$ gives the desired estimate.
\end{enumerate}

\textbf{Intuition:}
Each zero $\alpha$ of a function $f$ lets us "divide out" a factor $(z-\alpha)$ that is small. The Blaschke factor is the "perfect" way to divide out a zero in a disk while keeping the boundary value at 1. The more zeros you have, the "smaller" the function must be in the interior.
\end{solution}

\begin{exercise}
\label{analysis:2017:10}
Show that a single-valued analytic branch of $\sqrt{1-z^2}$ can be defined in any region such that the points $1$ and $-1$ are in the same component of the complement. What are the possible values of
\[
\int_\gamma \frac{dz}{\sqrt{1-z^2}}
\]
over a closed curve $\gamma$ in the region?
\end{exercise}

\begin{solution}
\textbf{Prerequisites:}
\begin{itemize}
    \item \textbf{Monodromy Theorem:} Conditions for defining single-valued branches of functions.
    \item \textbf{Residue Theorem:} Evaluating integrals based on the "winding" around poles.
\end{itemize}

\textbf{Steps:}
\begin{enumerate}
    \item \textbf{Branch criterion:} Write
    \[
    1-z^2=(1-z)(1+z).
    \]
    Along a closed curve $\gamma$ in the region, analytic continuation of the square root changes sign by
    \[
    (-1)^{\operatorname{Ind}(\gamma,1)+\operatorname{Ind}(\gamma,-1)}.
    \]
    Since $1$ and $-1$ lie in the same component of the complement, the winding number of $\gamma$ is constant on that component. Thus
    \[
    \operatorname{Ind}(\gamma,1)=\operatorname{Ind}(\gamma,-1),
    \]
    so the exponent is even and the continuation returns to the original value. Hence a single-valued analytic branch exists.
    \item \textbf{Periods of the differential:} On the model domain $\mathbb{C}\setminus[-1,1]$, choose a branch. The differential $dz/\sqrt{1-z^2}$ has only one possible period generator, represented by a loop going once around the cut $[-1,1]$.
    \item \textbf{Compute the generator:} Since
    \[
    \frac{d}{dz}\arcsin z=\frac{1}{\sqrt{1-z^2}}
    \]
    locally, analytic continuation once around both branch points changes the local inverse sine by $2\pi$ up to sign. Equivalently, integrating around a large circle and using the asymptotic
    \[
    \frac{1}{\sqrt{1-z^2}}\sim \frac{1}{\pm iz}
    \]
    gives a period $\pm 2\pi$. Therefore every closed curve has integral
    \[
    2\pi k,\qquad k\in\mathbb{Z},
    \]
    after choosing the sign of the branch; changing the branch only changes $k$ to $-k$.
\end{enumerate}

\textbf{Intuition:}
Taking a square root is like unrolling a double-covered surface. If you walk around \textbf{one} branch point, you end up on the "other side" of the surface. But if you walk around \textbf{both}, the two flips cancel out, and you end up back where you started.
\end{solution}

\begin{exercise}
\label{analysis:2017:11}
In $\mathbb{R}^n$, consider the Laplace equation
\[
u_{11}+u_{22}+\cdots+u_{nn}=0.
\]
Show that the equation is invariant under orthogonal transformations. Find all rotationally symmetric solutions to this equation.
\end{exercise}

\begin{solution}
\textbf{Prerequisites:}
\begin{itemize}
    \item \textbf{Laplacian in Polar Coordinates:} $\Delta u = u_{rr} + \frac{n-1}{r} u_r + \frac{1}{r^2} \Delta_{S} u$.
    \item \textbf{Orthogonal Transformation:} A linear map that preserves dot products and the gradient operator.
\end{itemize}

\textbf{Steps:}
\begin{enumerate}
    \item \textbf{Invariance:} Let $A$ be an orthogonal matrix and set $v(x)=u(Ax)$. By the chain rule,
    \[
    \nabla v(x)=A^T(\nabla u)(Ax).
    \]
    Differentiating once more gives
    \[
    D^2v(x)=A^T(D^2u)(Ax)A.
    \]
    Taking traces and using $AA^T=I$,
    \[
    \Delta v(x)=\operatorname{tr}(A^T(D^2u)(Ax)A)
    =\operatorname{tr}((D^2u)(Ax))=(\Delta u)(Ax).
    \]
    Thus $\Delta u=0$ if and only if $\Delta v=0$.
    \item \textbf{Radial Solutions:} For $u(r)$, the equation reduces to $u'' + \frac{n-1}{r} u' = 0$.
    \item \textbf{Solve ODE:} This is a separable ODE: $u''/u' = -(n-1)/r$. Integrating gives $\ln u' = -(n-1)\ln r + C$, so $u' = A r^{1-n}$.
    \item \textbf{Integrate again:}
    \begin{itemize}
        \item If $n=1$, $u(r)=Ar+B$ on each interval away from $0$.
        \item If $n=2$, $u(r) = A \ln r + B$.
        \item If $n \geq 3$, $u(r) = A r^{2-n} + B$.
    \end{itemize}
\end{enumerate}

\textbf{Intuition:}
The Laplace equation describes equilibrium states (like steady temperature). If the world is the same in all directions (rotationally symmetric), then the equilibrium should only depend on the distance from the center. In 2D, the potential grows slowly (log), while in 3D and above, it decays like a power of the distance.
\end{solution}

\begin{exercise}
\label{analysis:2017:12}
Let $A$ be an operator on $C[a,b]$ defined by
\[
(Af)(x)=\int_a^b K(x,y)f(y)\,dy,
\]
where $K(x,y)$ is continuous on $[a,b]\times[a,b]$ except at finitely many curves
\[
y=\phi_k(x),\qquad x\in[a,b],\qquad k=1,\ldots,n,
\]
with each $\phi_k$ continuous. Show that $A$ is a compact operator.
\end{exercise}

\begin{solution}
\textbf{Prerequisites:}
\begin{itemize}
    \item \textbf{Arzelà-Ascoli Theorem:} A set of functions is compact if it is bounded and equicontinuous.
    \item \textbf{Hilbert-Schmidt Operator:} An operator with a square-integrable kernel is compact on $L^2$.
\end{itemize}

\textbf{Steps:}
\begin{enumerate}
    \item \textbf{Boundedness convention:} The statement is understood in the usual jump-discontinuity sense: $K$ is bounded on the compact square and has only finite discontinuities along the listed curves. Then $A$ is a bounded map from $C[a,b]$ to $C[a,b]$ once continuity of $Af$ is checked.
    \item \textbf{$L^1$ continuity in the first variable:} If $x_j\to x$, then for every $y$ not equal to one of the finitely many values $\phi_k(x)$, the point $(x,y)$ is a continuity point of $K$, and hence
    \[
    K(x_j,y)\to K(x,y).
    \]
    Since $K$ is bounded and the exceptional set in $y$ has measure zero, dominated convergence gives
    \[
    \int_a^b |K(x_j,y)-K(x,y)|\,dy\to0.
    \]
    \item \textbf{Equicontinuity of the image of a bounded set:} If $\|f\|_\infty\leq 1$, then
    \[
    |Af(x_j)-Af(x)|
    \leq \int_a^b |K(x_j,y)-K(x,y)|\,dy\to0.
    \]
    The convergence is independent of the particular $f$ in the unit ball. Thus $A$ maps the unit ball of $C[a,b]$ to an equicontinuous family.
    \item \textbf{Uniform boundedness:} Also,
    \[
    |Af(x)|\leq (b-a)\|K\|_\infty\|f\|_\infty.
    \]
    Therefore the image of the unit ball is uniformly bounded and equicontinuous. By Arzel\`a--Ascoli, its closure in $C[a,b]$ is compact, so $A$ is compact.
\end{enumerate}

\textbf{Intuition:}
An integral operator is a "weighted average." If the weights (the kernel) are mostly smooth, the averaging process smooths out any input function. This smoothing turns a messy, infinite-dimensional set of functions into a much "smaller," more organized set, which is the definition of a compact operator.
\end{solution}
